\documentclass[11pt]{article}
\usepackage[T1]{fontenc}
\usepackage{lmodern}
\usepackage{microtype}
\usepackage[a4paper,margin=2.3cm]{geometry}
\usepackage{amsmath,amssymb}
\usepackage{graphicx}
\usepackage{booktabs}
\usepackage{enumitem}
\usepackage{xcolor}
\usepackage{caption}
\usepackage[most]{tcolorbox}
\usepackage{fancyhdr}
\usepackage{titlesec}
\usepackage[colorlinks=true,linkcolor=azul,citecolor=azul,urlcolor=azul]{hyperref}
\definecolor{azul}{HTML}{1D4ED8}
\definecolor{tinta}{HTML}{14213D}
\definecolor{ambar}{HTML}{B45309}
\definecolor{gris}{HTML}{64748B}
\definecolor{fondo}{HTML}{EEF2F7}
\graphicspath{{../../02_RESULTADOS/Graficos/presentacion/}{../../02_RESULTADOS/Animaciones_MP4/}}
\renewcommand{\figurename}{Figura}
\renewcommand{\tablename}{Tabla}
\renewcommand{\contentsname}{Índice}
\renewcommand{\refname}{Referencias}
\renewcommand{\appendixname}{Anexo}
\titleformat{\section}{\Large\bfseries\color{tinta}}{\thesection}{0.8em}{}
\titleformat{\subsection}{\large\bfseries\color{tinta}}{\thesubsection}{0.8em}{}
\setlist{itemsep=2pt,topsep=4pt}
\captionsetup{font=small,labelfont=bf}
\newtcolorbox{idea}[1][]{colback=fondo,colframe=fondo,coltitle=tinta,fonttitle=\bfseries,title=#1,sharp corners,boxrule=0pt,left=6pt,right=6pt,top=4pt,bottom=4pt}
\newtcolorbox{alerta}[1][]{colback=ambar!8,colframe=ambar,fonttitle=\bfseries,title=#1,sharp corners,boxrule=0.6pt,left=6pt,right=6pt}
\newcommand{\md}{\dot m}
\newcommand{\Qgw}{\dot Q_{gw}}
\newcommand{\Qwa}{\dot Q_{wa}}
\newcommand{\cvi}{c_v^{ig}}
\newcommand{\dd}{\mathrm{d}}
\pagestyle{fancy}
\fancyhf{}
\renewcommand{\headrulewidth}{0pt}
\fancyfoot[L]{\small\color{gris}ICF130 · Tema 7 · material interno del grupo}
\fancyfoot[R]{\small\color{gris}\thepage}
\setlength{\parindent}{0pt}
\setlength{\parskip}{5pt}

\title{\color{tinta}\textbf{Llenado transiente de un depósito con gas real}\\[4pt]
\large Informe técnico interno · ICF130 Tema 7 · material de estudio del grupo}
\author{Paul Berdichewsky · Benjamín Arancibia · Andrés Cerón\\ \small UTFSM, Campus Santiago San Joaquín · 2026-2}
\date{Noviembre 2026}
\begin{document}
\maketitle
\begin{alerta}[Uso de este documento]
Material privado para estudiar y defender el proyecto. No es la entrega. Todas las cifras salen de la
ejecución de \texttt{main.py} registrada en \texttt{02\_RESULTADOS/registro\_ejecucion.json}
(Python 3.13, NumPy 2.3.3, SciPy 1.16.1, CoolProp 7.2.0). Las derivaciones completas también están en
\texttt{Desarrollo\_Matematico.pdf}; las preguntas de defensa, en \texttt{Preguntas\_Profesor.pdf}.
\end{alerta}
\tableofcontents
\newpage

\section{Introducción}
\subsection{Problema}
Un depósito rígido de 50 L con nitrógeno a 1 bar y 300 K se conecta a un suministro a 200 bar y 300 K. El gas
entra, la presión sube y, aunque el gas que entra está a temperatura ambiente, el contenido se calienta.
Queremos saber cuánto demora en llegar a una presión objetivo, en qué estado queda (temperatura, masa,
factor $Z$) y cuánto pesan en eso la ecuación de estado, la restricción del caudal, la pared y el suministro.

\subsection{Importancia física}
Es el problema de llenar cilindros de gas comprimido: GNC, hidrógeno vehicular, aire de buceo o nitrógeno
industrial. En la práctica el calentamiento limita la velocidad de carga (los estándares de hidrógeno fijan
un máximo de 85 °C en el gas) y reduce cuánta masa cabe a una presión dada, porque un gas caliente necesita
menos masa para alcanzar esa presión.

\subsection{Objetivos}
\begin{enumerate}
\item Formular los balances transientes de masa y energía de un volumen de control abierto (exigencia del Tema 7).
\item Comparar gas ideal, Van der Waals y Peng-Robinson con una energía interna consistente con cada EOS.
\item Cuantificar el efecto de la restricción del caudal (V2) y del intercambio con una pared (V3).
\item Estudiar el efecto del estado del suministro $(P_s,T_s)$ (exigencia del Tema 7).
\item Verificar la implementación contra casos conocidos y estimar el orden de convergencia (exigencias transversales).
\end{enumerate}

\subsection{Por qué versiones sucesivas}
Cada versión agrega un solo fenómeno y mantiene todo lo demás igual. Así una diferencia entre versiones
se puede atribuir a una hipótesis. Dentro de cada versión, las tres EOS usan el mismo suministro, objetivo,
conexión y $\cvi$, para que la diferencia entre ellas venga del estado del gas.

\section{Sistema, hipótesis y dimensionalidad}
\begin{table}[h]\centering\small
\begin{tabular}{lll}\toprule
Magnitud & Valor & Origen\\\midrule
Volumen $V$ & 0,050 m$^3$ & condición de estudio\\
Estado inicial & 1 bar, 300 K (N$_2$ purgado) & condición de estudio\\
Suministro & 200 bar, 300 K, constante & condición de estudio\\
Objetivo & V1: 100 bar · V2, V3: 190 bar · malla: 140 bar & elección del grupo\\
$c_p$, $\cvi$, $\gamma$ & 1039 · 742,197 J/(kg K) · 1,39990 & NIST Shomate a 300 K\\
$R=R_u/M$ & 296,803 J/(kg K) & CODATA 2022\\
$T_c$, $P_c$, $\omega$ & 126,192 K · 3,3958 MPa · 0,0372 & Span et al. (2000)\\
Caudal V1 & 20 g/s & elección\\
Conexión V2-V3 & $C_d=0{,}8$, $A=5{,}447\cdot10^{-7}$ m$^2$ (20 g/s estrangulado) & calibración\\
Pared V3 & 8 kg, $c_w=500$ J/(kg K), $C_w=4000$ J/K & Outokumpu 304 (solo $c_w$)\\
Conductancias & $UA_{gw}=8$ W/K, $UA_{wa}=3$ W/K, $T_\infty=300$ K & parámetros de estudio\\\bottomrule
\end{tabular}
\caption{Caso base. Masa de pared, conductancias y $C_dA$ no son medidas: por eso se estudian con sensibilidad.}
\end{table}

\textbf{Hipótesis.} Depósito rígido; una entrada y ninguna salida; sin trabajo de eje; gas bien mezclado;
suministro de remanso constante (depósito infinito); energía cinética y potencial almacenadas despreciables;
N$_2$ puro, monofásico y muy por encima de $T_c$; $\cvi$ constante como fondo ideal de las tres EOS;
en V3, una sola temperatura de pared.

\subsection{Modelo 0D y visualización}
Todo el proyecto es un modelo \textbf{0D} (parámetros concentrados): las variables dependen del tiempo pero no
de la posición. No hay partes 1D ni 2D. La animación pinta cada estanque con un color \emph{uniforme} igual a
$T_g(t)$ y el anillo con $T_w(t)$: representa exactamente lo que se calcula. La animación del avance usaba un
mapa con un punto caliente en la entrada; se retiró porque (i) no proviene de ninguna ecuación resuelta y
(ii) contradice la física: el gas que entra es el más frío (300 K o menos tras la expansión), y el
calentamiento se debe a la compresión del gas residente. Un campo espacial exigiría resolver Navier-Stokes y
energía (CFD) o, al menos, un modelo 1D de conducción en la pared, que no forman parte del proyecto.
La referencia a von Kármán que aparecía en el avance (calle de vórtices tras una esfera) no tenía uso en el
modelo ni en el código y se eliminó: una calle de vórtices requiere un cuerpo en un flujo cruzado.

\section{Fundamentos y desarrollo matemático}
% Desarrollo matemático completo. Se incluye en el informe y en el documento aparte.

\subsection{Notación y unidades}
Todas las magnitudes están en SI. $m$ masa de gas [kg]; $V$ volumen del depósito [m$^3$];
$v=V/m$ volumen específico [m$^3$/kg]; $T_g$ temperatura del gas [K]; $T_w$ temperatura media de la
pared [K]; $P$ presión absoluta [Pa]; $u$ y $h$ energía interna y entalpía específicas [J/kg];
$U=mu$ energía interna del gas [J]; $\md$ caudal que entra [kg/s]; $\Qgw$ potencia térmica del gas a la
pared y $\Qwa$ de la pared al ambiente [W], ambas positivas en el sentido gas $\to$ pared $\to$ ambiente.
El subíndice $s$ indica el suministro (remanso) y el $0$ el estado inicial. $R=R_u/M=296{,}803$ J/(kg K).

\subsection{Balance de masa}
Para un volumen de control con entradas y salidas, $\dd m/\dd t=\sum \md_{\text{ent}}-\sum \md_{\text{sal}}$.
Hay una sola entrada y ninguna salida:
\begin{equation}\label{eq:masa}
\frac{\dd m}{\dd t}=\md .
\end{equation}

\subsection{Primera ley para el volumen de control}
La forma general, despreciando la energía cinética y potencial \emph{almacenadas} en el gas, es
\[
\frac{\dd U}{\dd t}=\dot Q_{\text{neto}}-\dot W_{\text{eje}}-\dot W_{\text{frontera}}
+\sum \md_{\text{ent}}\Big(h+\tfrac{c^2}{2}+gz\Big)_{\text{ent}}-\sum \md_{\text{sal}}\Big(h+\tfrac{c^2}{2}+gz\Big)_{\text{sal}} .
\]
\begin{itemize}
\item $\dot W_{\text{eje}}=0$: no hay compresor, turbina ni agitador dentro del depósito.
\item $\dot W_{\text{frontera}}=-P\,\dd V/\dd t=0$: el depósito es rígido.
\item La corriente se toma en su estado de \emph{remanso} $(P_s,T_s)$: la entalpía de remanso $h_s=h+c^2/2$
ya incluye la energía cinética del chorro, que se disipa dentro del depósito. $gz$ es despreciable.
\item $\dot Q_{\text{neto}}=-\Qgw$ (el gas pierde calor hacia la pared en V3; vale cero en V1 y V2).
\end{itemize}
Queda
\begin{equation}\label{eq:energia}
\frac{\dd (mu)}{\dd t}=\md\,h_s-\Qgw .
\end{equation}

\begin{idea}[¿Por qué aparece la entalpía y no la energía interna?]
Para meter un kilo de gas al volumen de control hay que empujarlo contra la presión: el gas que viene
detrás hace un trabajo $P_s v_s$ por kilo (trabajo de flujo). La energía que cruza la frontera por kilo es
$u_s+P_sv_s=h_s$. Ese trabajo de flujo es el que calienta el depósito aunque no entre calor:
en un gas ideal $h_s-u_s=RT_s$, que es la diferencia entre $c_pT_s$ y $c_vT_s$.
Energía interna = energía almacenada; entalpía = energía transportada por la corriente;
calor = energía que cruza la frontera por diferencia de temperatura, sin masa.
\end{idea}

\subsection{Pared y ambiente (V3)}
La pared se representa con una capacidad térmica $C_w=m_wc_w=8\cdot500=4000$ J/K y una sola temperatura.
Las transferencias se escriben con conductancias globales $UA$ [W/K]:
\begin{equation}\label{eq:pared}
\Qgw=UA_{gw}(T_g-T_w),\qquad \Qwa=UA_{wa}(T_w-T_\infty),\qquad C_w\frac{\dd T_w}{\dd t}=\Qgw-\Qwa .
\end{equation}
Sumando \eqref{eq:energia} y \eqref{eq:pared} se cancela $\Qgw$ (transferencia interna):
$\dd(U_g+U_w)/\dd t=\md h_s-\Qwa$. Esa identidad se usa como prueba del balance global.

\subsection{Relación entre la ecuación de estado y la energía interna}
De $\dd u=T\dd s-P\dd v$ y de la relación de Maxwell $(\partial s/\partial v)_T=(\partial P/\partial T)_v$:
\begin{equation}\label{eq:maxwell}
\dd u=c_v\,\dd T+\Big[T\Big(\frac{\partial P}{\partial T}\Big)_v-P\Big]\dd v ,
\qquad u_v\equiv\Big(\frac{\partial u}{\partial v}\Big)_T=T\Big(\frac{\partial P}{\partial T}\Big)_v-P .
\end{equation}
Consecuencia: si se cambia la ecuación de estado, $u$ cambia. Usar una EOS real con $u=c_vT$ sería
inconsistente. Se integra $u_v$ a $T$ fija desde el gas diluido ($v\to\infty$, donde $u=\cvi T$) hasta $v$:
$u(T,v)=\cvi T+\int_\infty^{v}u_v\,\dd v'$.

\paragraph{Gas ideal.} $P=RT/v$, $T(\partial P/\partial T)_v=P$, $u_v=0$: $u=\cvi T$, $h=\cvi T+RT=c_pT$.

\paragraph{Van der Waals.} $P=\dfrac{RT}{v-b}-\dfrac{a}{v^2}$. Entonces $T(\partial P/\partial T)_v=RT/(v-b)$ y
$u_v=a/v^2$. Integrando, $u=\cvi T-a/v$. Las constantes salen de imponer
$(\partial P/\partial v)_T=(\partial^2P/\partial v^2)_T=0$ en el punto crítico:
\[
a=\frac{27R^2T_c^2}{64P_c}=174{,}28\ \text{Pa m}^6/\text{kg}^2,\qquad b=\frac{RT_c}{8P_c}=1{,}3787\cdot10^{-3}\ \text{m}^3/\text{kg}.
\]
En $U=mu=m\cvi T-am^2/V$ aparece el término $-am^2/V$.

\paragraph{Peng-Robinson (base másica).}
\begin{equation}\label{eq:pr}
P=\frac{RT}{v-b}-\frac{a(T)}{D},\quad D=v^2+2bv-b^2,\quad
a(T)=0{,}45724\frac{R^2T_c^2}{P_c}\,g^2,\quad g=1+\kappa\big(1-\sqrt{T/T_c}\big),
\end{equation}
con $b=0{,}07780\,RT_c/P_c$ y $\kappa=0{,}37464+1{,}54226\,\omega-0{,}26992\,\omega^2$, $\omega=0{,}0372$.
Derivadas: $g'=-\kappa/(2\sqrt{TT_c})$, $g''=-g'/(2T)$, $a'=2a_cgg'$, $a''=2a_c(g'^2+gg'')$.
Como $(\partial P/\partial T)_v=R/(v-b)-a'/D$:
\[
u_v=T\frac{R}{v-b}-\frac{Ta'}{D}-\frac{RT}{v-b}+\frac{a}{D}=\frac{a-Ta'}{D}.
\]
Con $I(v)=\dfrac{1}{2\sqrt2\,b}\ln\dfrac{v+(1+\sqrt2)b}{v+(1-\sqrt2)b}$ se cumple $I'(v)=-1/D$ e $I\to0$ cuando $v\to\infty$. Por lo tanto
\begin{equation}\label{eq:upr}
u=\cvi T+(Ta'-a)\,I(v),\qquad h=u+Pv,\qquad c_v=\Big(\frac{\partial u}{\partial T}\Big)_v=\cvi+Ta''\,I(v).
\end{equation}
Comprobación: $\partial u/\partial v=(Ta'-a)(-1/D)=(a-Ta')/D=u_v$. Si $b\to0$, $I\to1/v$ y se recupera la forma de Van der Waals.
En PR, $c_v$ depende del estado aunque $\cvi$ sea constante.

\subsection{Ecuación diferencial de la temperatura del gas}
Con $U=m\,u(T_g,v)$ y $v=V/m$, por la regla de la cadena:
\[
\frac{\dd U}{\dd t}=u\frac{\dd m}{\dd t}+m\Big(c_v\frac{\dd T_g}{\dd t}+u_v\frac{\dd v}{\dd t}\Big),\qquad
\frac{\dd v}{\dd t}=-\frac{V}{m^2}\frac{\dd m}{\dd t}=-\frac{v}{m}\,\md .
\]
Igualando con \eqref{eq:energia}:
\begin{equation}\label{eq:dTdt}
\boxed{\ \frac{\dd T_g}{\dd t}=\frac{\md\,\big[h_s-u+v\,u_v\big]-\Qgw}{m\,c_v}\ }
\end{equation}
\begin{itemize}
\item Gas ideal: $u_v=0$, $u=\cvi T_g$, $h_s=c_pT_s$ $\Rightarrow$ $\dfrac{\dd T_g}{\dd t}=\dfrac{\md(c_pT_s-\cvi T_g)-\Qgw}{m\cvi}$.
\item Van der Waals: $v\,u_v=a/v$, $u=\cvi T_g-a/v$ $\Rightarrow$ $h_s-u+vu_v=h_s-\cvi T_g+2a/v=h_s-\cvi T_g+2am/V$
(el término $2am/V$ del avance).
\item Con $a=b=0$ se recupera exactamente el gas ideal; con $UA_{gw}=UA_{wa}=0$, V3 se reduce a V2.
\end{itemize}
Las variables integradas son $(m,T_g,T_w)$. Como $m$ y $T_g$ se conocen en cada paso, $P$ se evalúa
directamente con la EOS, sin resolver ninguna ecuación implícita durante la integración.

\subsection{Solución exacta de V1 con gas ideal}
Con $\Qgw=0$ y $u=\cvi T$: $\dd(mT)/\dd m=\gamma T_s$, con $\gamma=c_p/\cvi$, pues
$\cvi\,\dd(mT)=c_pT_s\,\dd m$. Integrando desde $(m_0,T_0)$:
\begin{equation}\label{eq:v1}
T_g(m)=\gamma T_s+\frac{m_0}{m}\big(T_0-\gamma T_s\big),\qquad m=m_0+\md t .
\end{equation}
Cuando $m\gg m_0$, $T_g\to\gamma T_s=1{,}39990\cdot300=419{,}97$ K. Además $PV=mRT_g$ da
$P(t)=P_0+R\gamma T_s\md\,t/V$: la presión crece linealmente. El tiempo exacto para 100 bar es
$t_f=(P_f-P_0)V/(R\gamma T_s\md)=198{,}559$ s, igual al numérico hasta $2\cdot10^{-12}$ s.
En V2 la misma $T_g(m)$ sigue valiendo (el caudal solo cambia el reloj) y el tiempo se obtiene con
$t_f=\int_{P_0}^{P_f}\dd P/\big[(R\gamma T_s/V)\,\md(P)\big]$.

\subsection{Caudal compresible (V2 y V3)}
Se modela la conexión como una tobera isentrópica de gas ideal entre el remanso $(P_s,T_s)$ y la presión
del depósito $P$. Con $r=P/P_s$, la temperatura en la sección de salida es $T_e=T_sr^{(\gamma-1)/\gamma}$; la
energía da $c_e^2/2=c_p(T_s-T_e)$ y la continuidad $\md=C_dA\,\rho_ec_e$ con $\rho_e=P/(RT_e)$. Sustituyendo:
\begin{equation}\label{eq:caudal}
\md=C_dA\frac{P_s}{\sqrt{RT_s}}\sqrt{\frac{2\gamma}{\gamma-1}\Big(r^{2/\gamma}-r^{(\gamma+1)/\gamma}\Big)}\quad (r^*<r<1).
\end{equation}
Esta expresión tiene un máximo en $r^*=\big(2/(\gamma+1)\big)^{\gamma/(\gamma-1)}=0{,}5283$, donde la velocidad
es sónica. Para $r\le r^*$ el caudal queda fijo (flujo estrangulado) en
\[
\md^*=C_dA\,P_s\sqrt{\frac{\gamma}{RT_s}}\,\Big(\frac{2}{\gamma+1}\Big)^{\frac{\gamma+1}{2(\gamma-1)}} :
\]
bajar más la presión del depósito no aumenta el caudal. Las dos ramas empalman con continuidad de valor y
de pendiente. $C_dA=4{,}358\cdot10^{-7}$ m$^2$ se calibra para $\md^*=20$ g/s en el caso base y no se cambia
nunca. El cambio de régimen ocurre en $P=r^*P_s=105{,}66$ bar.

\subsection{Factor de compresibilidad y masa a presión fija}
$Z=Pv/(RT)$. A igual $P$ y $V$: $m=PV/(ZRT)$. Dos modelos que llegan a la misma presión con distinta
$Z$ y $T$ guardan distinta masa; la comparación depende del producto $ZT$, no solo de $T$.

\subsection{Límite isotermo}
Si un baño perfecto mantiene $T_g=T_0$, la EOS da $m(P)=V/v(P,T_0)$ y
$\dd t=\dd m/\md=-\dfrac{V}{v^2(\partial P/\partial v)_T\,\md(P)}\,\dd P$.
El calor que hay que extraer sale de la primera ley integrada:
$Q_{\text{ext}}=(m_f-m_0)h_s-\big[U(m_f)-U(m_0)\big]$. En gas ideal con $T_s=T_0$:
$Q_{\text{ext}}=(m_f-m_0)(c_p-\cvi)T_0=(m_f-m_0)RT_0$, exactamente el trabajo de flujo acumulado.

\subsection{Residuos de los balances}
Con las tasas reconstruidas de la solución e integradas con cuadratura de Gauss-Legendre (8 nodos por paso):
\begin{align*}
R_m&=\Delta m-\!\int\!\md\,\dd t, & R_g&=\Delta U_g-\!\int\!\md h_s\,\dd t+Q_{gw},\\
R_w&=\Delta U_w-Q_{gw}+Q_{wa}, & R_\Sigma&=\Delta U_g+\Delta U_w-\!\int\!\md h_s\,\dd t+Q_{wa}.
\end{align*}
Se normalizan con la masa añadida o con la suma de los módulos de los términos de cada balance.
Un residuo pequeño prueba que la ecuación se resolvió bien, no que la EOS sea correcta.

\subsection{Diferencias porcentuales}
Cuando se reporta una diferencia porcentual entre un modelo $X$ y una referencia $Y$ se usa
$\Delta\%=100\,(X-Y)/Y$. En el contraste con Span la referencia es siempre Span; en la malla de suministro, el
modelo de referencia del par (ideal o Van der Waals). Las energías residuales se comparan en kJ/kg y no en
porcentaje porque cerca de cero el porcentaje no es informativo.


\section{Desarrollo de cada versión}
\subsection{V1: caudal fijo, adiabático, 100 bar}
\textbf{Qué se implementó.} $\md=20$ g/s, $\Qgw=0$; EDO \eqref{eq:dTdt} con las tres EOS.
\textbf{Para qué.} Aislar el efecto de la EOS: el caudal no depende de la presión.
\textbf{Resultado.} El gas ideal sube casi de inmediato hacia $\gamma T_s\approx420$ K, porque la masa inicial
(0,056 kg) es despreciable frente a la que entra. Las EOS reales terminan unos 30 K más frías (388-389 K) y
necesitan más masa (4,23-4,32 kg frente a 4,03 kg) y más tiempo (208-213 s frente a 199 s).
\textbf{Limitación.} Un caudal fijo hasta 100 bar no es físico si el suministro está a 200 bar con una conexión real.

\subsection{V2: caudal compresible, adiabático, 190 bar}
\textbf{Qué cambia.} $\md(P)$ según \eqref{eq:caudal}. El flujo es estrangulado hasta 105,66 bar (210 s en gas
ideal) y luego subsónico; al llegar a 190 bar el caudal cayó a 8,99 g/s en las tres EOS (a igual $P$, igual caudal).
\textbf{Resultado.} 418,92 s y 419,09 K (ideal); 410,86 s y 403,67 K (VdW); 404,88 s y 402,67 K (PR).
Ahora las EOS reales llegan \emph{antes} con \emph{menos} masa (sección~\ref{sec:porque}).
\textbf{Aporte.} El caudal depende del estado del depósito: la masa deja de crecer linealmente.

\subsection{V3: pared con inercia térmica y ambiente, 190 bar}
\textbf{Qué cambia.} Se agregan $T_w$ y \eqref{eq:pared}. \textbf{Escalas de tiempo:} el gas responde en
$m_0c_v/(\md c_v+UA_{gw})\approx1{,}8$ s y la pared en $C_w/(UA_{gw}+UA_{wa})\approx364$ s.
\textbf{Resultado.} 456,40 s / 385,32 K (ideal); 440,75 s / 379,44 K (VdW); 433,60 s / 378,59 K (PR).
La pared termina en 343,5 K (ideal) y 334,6 K (PR). Respecto de V2 la temperatura final baja 33,8 K (ideal),
24,2 K (VdW) y 24,1 K (PR) y el llenado se alarga entre 29 y 37 s.
\textbf{Aporte.} El gas cede 208,6 kJ (ideal) o 161,8 kJ (PR) a la pared, que solo pierde 34,6 / 23,6 kJ al ambiente:
la pared actúa sobre todo como almacén. En gas ideal $T_g$ tiene un máximo de 385,97 K en 378 s y luego baja
0,65 K: cuando el caudal cae, el aporte de entalpía ya no compensa la pérdida a la pared.

\section{Implementación numérica}
\begin{itemize}
\item \textbf{Integrador:} \texttt{solve\_ivp} (SciPy) con RK45, par embebido de órdenes 5 y 4 con control de error
local $\le \text{atol}+\text{rtol}|y|$; rtol $=10^{-9}$, atol $=10^{-11}$, paso máximo 0,5 s (resuelve el transiente
inicial de 1,8 s y la transición de régimen).
\item \textbf{Evento:} función $P(t)-P_{obj}$ con \texttt{terminal=True} y \texttt{direction=1}: el instante se ubica sobre
la solución densa (no en la malla de salida). Con malla de 10 s y umbral el error sería de hasta 9,25 s; con el evento es despreciable.
\item \textbf{Raíces:} los estados definidos por $(P,T)$ (inicial y suministro) se resuelven en densidad con Brent en
$\rho\in[0,(1-10^{-10})/b)$, donde la función cambia de signo. Se exige $v>b$ y $(\partial P/\partial v)_T<0$ (rama estable);
como $T>T_c$ la raíz es única. En PR también se exige $c_v>0$ en cada evaluación.
\item \textbf{Formulación alternativa (PR):} se integra $(m,U,T_w)$ con DOP853 y se recupera $T$ resolviendo $m\,u(T,V/m)=U$.
Coincide con la formulación en temperatura hasta $1{,}6\cdot10^{-6}$ K y $3\cdot10^{-9}$ s.
\item \textbf{Orden de convergencia:} RK4 clásico de paso fijo en V3 ideal hasta 180 s (antes de la transición),
contra DOP853 con rtol $2\cdot10^{-13}$. Si $e(h)\approx Kh^p$, $p=\log_2[e(h)/e(h/2)]$: se obtuvo 4,14 a 4,03.
\item \textbf{Estabilidad:} el sistema no es rígido en la región estudiada (relación de escalas $\sim$200); los métodos
explícitos funcionaron. La conclusión vale para estos parámetros.
\end{itemize}

\subsection{Estructura del código}
\texttt{main.py} es el único punto de entrada. Módulos:
\begin{itemize}
\item \texttt{modelos/base.py}: ideal y VdW, copia sin cambios físicos de la entrega anterior;
\item \texttt{modelos/peng\_robinson.py}: PR, Span, malla de suministro, isotermo;
\item \texttt{modelos/sensibilidad\_pr.py} y \texttt{modelos/figuras\_presentacion.py};
\item \texttt{modelos/animacion.py}: video MP4;
\item \texttt{pruebas/test\_modelo.py}: 25 pruebas automáticas.
\end{itemize}
Ejecución completa: 63 s sin video y 92 s el video, en el entorno de prueba.

\section{Resultados}
\begin{table}[h]\centering\small
\begin{tabular}{llrrrrr}\toprule
Etapa & EOS & $P_{obj}$ [bar] & $t_f$ [s] & $T_f$ [K] & $m_f$ [kg] & $Z_f$\\\midrule
V1 & Ideal & 100 & 198,56 & 418,30 & 4,027 & 1,0000\\
 & Van der Waals & 100 & 213,36 & 387,92 & 4,323 & 1,0045\\
 & Peng-Robinson & 100 & 208,53 & 389,49 & 4,227 & 1,0233\\
V2 & Ideal & 190 & 418,92 & 419,09 & 7,637 & 1,0000\\
 & Van der Waals & 190 & 410,86 & 403,67 & 7,595 & 1,0439\\
 & Peng-Robinson & 190 & 404,88 & 402,67 & 7,471 & 1,0640\\
V3 & Ideal & 190 & 456,40 & 385,32 & 8,307 & 1,0000\\
 & Van der Waals & 190 & 440,75 & 379,44 & 8,132 & 1,0374\\
 & Peng-Robinson & 190 & 433,60 & 378,59 & 7,985 & 1,0587\\\bottomrule
\end{tabular}
\caption{Estado al alcanzar la presión objetivo (\texttt{Datos/peng\_robinson/casos\_base.csv}).}
\end{table}

\begin{figure}[h]\centering
\includegraphics[width=0.49\linewidth]{f04_V1_T_P.png}\hfill
\includegraphics[width=0.49\linewidth]{f03_V2_caudal.png}
\caption{Izquierda: V1, temperatura y presión. Derecha: V2, caudal y cambio de régimen.}
\end{figure}
\begin{figure}[h]\centering
\includegraphics[width=0.49\linewidth]{f01_V3_temperaturas.png}\hfill
\includegraphics[width=0.49\linewidth]{f02_V3_presion_masa.png}
\caption{V3: temperaturas de gas y pared (izquierda); presión y masa (derecha).}
\end{figure}

\textbf{Lectura.} En V1 la presión del gas ideal crece linealmente (solución exacta) y en las EOS reales se curva.
En V2 la masa crece linealmente mientras el flujo está estrangulado y luego se curva. En V3 la pared retrasa
el calentamiento: el gas queda en una meseta y la pared sube casi linealmente.

\section{Comparaciones}
\subsection{Por qué difieren las ecuaciones de estado}\label{sec:porque}
\begin{enumerate}
\item \textbf{Entalpía del suministro.} A 200 bar y 300 K, $h_s$ vale 311,7 kJ/kg (ideal), 273,5 (VdW) y 275,4 (PR)
con la misma referencia $\cvi T$. La atracción entre moléculas da una entalpía residual negativa
(Span: $-32{,}3$ kJ/kg). Entra menos energía por kilo: las EOS reales terminan más frías.
\item \textbf{Compresibilidad.} $m=PV/(ZRT)$. A baja presión domina la atracción ($Z<1$ en VdW) y a 190 bar el volumen
excluido ($Z=1{,}037$ VdW, $1{,}059$ PR). En V1 (100 bar) el efecto de la menor $T$ domina y las EOS reales
necesitan más masa; en V2-V3 (190 bar) domina $Z$ y necesitan menos, por eso llegan antes.
\end{enumerate}
\begin{figure}[h]\centering
\includegraphics[width=0.49\linewidth]{f06_V3_Z.png}\hfill
\includegraphics[width=0.49\linewidth]{f05_V2_vs_V3.png}
\caption{Factor $Z$ en V3 (izquierda) y efecto de la pared: V2 punteado, V3 continuo (derecha).}
\end{figure}

\subsection{Sensibilidad a la pared (tres EOS, 190 bar)}
\begin{table}[h]\centering\small
\begin{tabular}{lrrrrrr}\toprule
Parámetro & valor & $t_f$ ideal [s] & $t_f$ PR [s] & $T_{max}$ ideal [K] & $T_{max}$ PR [K] & $Q_{gw}$ PR [kJ]\\\midrule
$UA_{gw}$ [W/K] & 0 & 418,92 & 404,88 & 419,09 & 402,67 & 0\\
 & 4 & 444,26 & 423,58 & 396,90 & 386,73 & 104,4\\
 & 8 & 456,40 & 433,60 & 385,97 & 378,59 & 161,8\\
 & 16 & 466,26 & 442,95 & 376,85 & 371,14 & 217,2\\
$UA_{wa}$ [W/K] & 1,5 & 455,54 & 433,07 & 386,39 & 378,95 & 159,3\\
 & 6 & 457,94 & 434,56 & 385,24 & 377,96 & 166,4\\
$C_w$ [J/K] & 2000 & 449,27 & 428,34 & 391,29 & 382,55 & 133,5\\
 & 8000 & 462,66 & 437,98 & 382,16 & 375,49 & 184,6\\\bottomrule
\end{tabular}
\caption{Un parámetro a la vez desde el caso base (\texttt{Datos/sensibilidad\_3EOS}).}
\end{table}
$UA_{gw}$ es el parámetro dominante. $UA_{wa}$ casi no influye en este horizonte (la pared todavía se está calentando:
el calor queda almacenado). Una pared de mayor $C_w$ se calienta menos, mantiene más diferencia con el gas y extrae más calor.
Las tres EOS siguen la misma tendencia.

\subsection{Extremos adiabático e isotermo}
Con un baño ideal a 300 K y el mismo caudal, el llenado a 190 bar tarda 586,4 s (ideal), 581,0 (VdW) y 569,5 (PR),
cerca de 40\,\% más que el adiabático, entra cerca de 40\,\% más masa y hay que extraer 945,0, 937,3 y 936,2 kJ.
En gas ideal se comprobó $Q=(m_f-m_0)RT_0$ hasta $10^{-8}$ J. V3 queda entre ambos extremos.

\subsection{Estado del suministro}
Malla $P_s\in\{150,175,200,225,250\}$ bar $\times$ $T_s\in\{280,\dots,320\}$ K, objetivo común 140 bar, misma $C_dA$.
El caudal estrangulado crece con $P_s$ y cae como $T_s^{-1/2}$; en las EOS reales $P_s$ y $T_s$ también cambian $h_s$.
PR $-$ ideal en tiempo va de $-15{,}5$ s (150 bar, 320 K) a $+8{,}4$ s (250 bar, 280 K): \textbf{el signo cambia}.
En temperatura final PR siempre termina más frío (de $-3{,}7$ a $-21{,}3$ K). Frente a VdW, PR es más rápido y
guarda menos masa en los 25 pares. A 200 bar y 300 K, cambiar solo el depósito a PR (con $h_s$ ideal) acorta
26,57 s y cambiar después $h_s$ a PR alarga 25,86 s: el neto $-0{,}72$ s esconde dos efectos grandes.
La descomposición depende del orden porque caudal, presión y energía están acoplados.
\begin{figure}[h]\centering
\includegraphics[width=0.45\linewidth]{f08_malla_tiempo.png}\hfill
\includegraphics[width=0.45\linewidth]{f10_descomposicion.png}
\caption{Diferencia PR $-$ ideal en el tiempo hasta 140 bar y descomposición a 200 bar, 300 K.}
\end{figure}

\subsection{Contraste con la ecuación de referencia de Span}
Se comparó con HEOS::Nitrogen (CoolProp 7.2.0) a igual $(T,P)$ en 35 estados y en los nueve estados finales.
A 200 bar y 300 K ($\rho_{ref}=212{,}54$ kg/m$^3$): error de densidad $+5{,}68$\,\% (ideal), $+4{,}90$\,\% (VdW),
$+2{,}83$\,\% (PR); error de $h$ residual $+32{,}3$, $-5{,}9$ y $-4{,}0$ kJ/kg. Máximos en la malla: densidad 12,7 / 5,5 / 3,7\,\%.
PR es la más cercana \emph{en este rango}; no es exacta.
\begin{figure}[h]\centering
\includegraphics[width=0.8\linewidth]{f11_contraste_span_300K.png}
\caption{Errores frente a Span a 300 K.}
\end{figure}

\section{Verificación}
\begin{itemize}
\item Solución exacta V1 ideal: error máximo $2{,}56\cdot10^{-6}$ K; $t_f$ con error $2\cdot10^{-12}$ s. V2 ideal: $T(m)$ y $t_f$ por cuadratura, error $2{,}3\cdot10^{-6}$ K.
\item Límites: $a=b=0$ recupera el ideal en V1-V3 (VdW y PR); V3 con $UA=0$ reproduce V2 (tres EOS).
\item Balances: residuo energético relativo máximo $8\cdot10^{-12}$ en los casos base y $<10^{-10}$ en los 28 casos PR (3 base y 25 de la malla).
\item Convergencia: barridos de paso máximo y tolerancia (RK45 y DOP853) contra referencia refinada; RK4 de orden 4,03-4,14.
\item Propiedades PR en 20 estados: $h=u+Pv$, cuadratura independiente de $u_v$, $c_v$ y $u_v$ por diferencias centradas.
\item Reproducibilidad: los 74 CSV regenerados coinciden con los guardados por el grupo con diferencias relativas $<10^{-6}$ en las
magnitudes físicas; solo cambian columnas de error a nivel de redondeo.
\end{itemize}
\begin{figure}[h]\centering
\includegraphics[width=0.45\linewidth]{f12_orden_rk4.png}\hfill
\includegraphics[width=0.5\linewidth]{f13_balances.png}
\caption{Orden de RK4 y cierre del balance de energía.}
\end{figure}

\section{Limitaciones y trabajo futuro}
\begin{itemize}
\item \textbf{0D:} no hay estratificación ni chorro; en estanques grandes o llenados rápidos la temperatura local
puede diferir del promedio. Futuro: modelo 1D de pared o CFD axisimétrico.
\item \textbf{Conexión:} ley isentrópica de gas ideal para las tres EOS (hipótesis de comparación). A 200 bar un caudal
con propiedades reales cambiaría unos puntos porcentuales; su efecto debería aislarse en otra comparación.
\item \textbf{Parámetros:} $UA$, $m_w$ y $C_dA$ no medidos. Se necesitan geometría y mediciones de $P(t)$, $\md(t)$, $T_g$ y $T_w$.
\item \textbf{Propiedades:} $\cvi$ constante (en 300-420 K el $c_p$ ideal del N$_2$ varía menos de 1\,\%); pared con $c_w$ constante.
\item \textbf{Suministro y válvula:} suministro infinito y apertura instantánea; sin agotamiento del suministro.
\item \textbf{Sin validación experimental:} las pruebas confirman que se resuelven bien las ecuaciones, no que el modelo prediga un estanque real.
\end{itemize}

\section{Conclusiones}
\begin{enumerate}
\item (Objetivo 1) El calentamiento sin calor externo se explica por el trabajo de flujo: la corriente aporta $h_s=u_s+P_sv_s$.
Sin pared y con gas ideal el gas tiende a $\gamma T_s\approx420$ K.
\item (Objetivo 2) Una EOS real exige una energía interna consistente (Maxwell). Con ella, a 190 bar las EOS reales dan
llenados más cortos y menos masa por dos mecanismos: menor entalpía de entrada y $Z>1$.
\item (Objetivo 3) La restricción del caudal domina la duración desde 106 bar; la pared baja la temperatura final 24-34 K
y alarga el llenado; su capacidad importa más que las pérdidas al ambiente.
\item (Objetivo 4) Según $(P_s,T_s)$ PR llega antes o después que el ideal; el resultado neto surge de efectos opuestos.
\item (Objetivo 5) La implementación reproduce soluciones exactas, límites y balances, converge con el orden teórico y,
frente a Span, PR es la EOS más cercana en el rango estudiado.
\end{enumerate}

\begin{thebibliography}{9}
\bibitem{tema7} S. Palma e I. Díaz. \emph{Proyecto computacional: descripción ampliada de los temas}, Tema 7. ICF130, UTFSM, 2026.
\bibitem{palma} S. Palma. \emph{Fundamentos de Termodinámica Aplicada}, material de clase (clases 11, 13-16). UTFSM, 2026.
\bibitem{pr} D.-Y. Peng y D. B. Robinson. A new two-constant equation of state. \emph{Ind. Eng. Chem. Fundam.} 15(1):59-64, 1976. doi:10.1021/i160057a011.
\bibitem{span} R. Span, E. W. Lemmon, R. T. Jacobsen, W. Wagner y A. Yokozeki. A reference equation of state for the thermodynamic properties of nitrogen. \emph{J. Phys. Chem. Ref. Data} 29(6):1361-1433, 2000. doi:10.1063/1.1349047.
\bibitem{nist} NIST Chemistry WebBook, SRD 69, Nitrogen (datos de fase gaseosa, Shomate). \url{https://webbook.nist.gov/cgi/cbook.cgi?ID=C7727379&Mask=1}. Consulta: 15-11-2026.
\bibitem{codata} NIST SP 959, CODATA 2022. \url{https://www.physics.nist.gov/cuu/pdf/wallet_2022.pdf}.
\bibitem{nasa} NASA Glenn Research Center. Mass flow rate equations. \url{https://www.grc.nasa.gov/WWW/k-12/airplane/mchkdrv.html}.
\bibitem{outo} Outokumpu. Core range datasheet, acero 304/4301. Consulta: 15-11-2026.
\bibitem{coolprop} I. H. Bell et al. CoolProp 7.2.0. \url{https://coolprop.org}. SciPy 1.16.1, \texttt{solve\_ivp} y \texttt{brentq}.
\end{thebibliography}

\appendix
\section{Ejecución y archivos}
\begin{verbatim}
cd 01_ENTREGA_PROFESOR/Codigo_Final
python -m pip install -r requirements.txt
python main.py                     # todo (~3 min)
python main.py --caso V3 --eos pr  # un caso, resumen en pantalla
python -m pytest                   # 25 pruebas
\end{verbatim}
Datos en \texttt{02\_RESULTADOS/Datos}, figuras en \texttt{Graficos}, video en \texttt{Animaciones\_MP4}.

\section{Cambios respecto del avance y de la entrega anterior}
Unificación en un solo \texttt{main.py}; Peng-Robinson completo y contraste con Span; malla Ps-Ts y descomposición;
límite isotermo; PR en la sensibilidad de pared; animación nueva con tres EOS y color uniforme (se retiró el mapa
ilustrativo y la referencia a von Kármán); 25 pruebas automáticas. La física de ideal y VdW no cambió.
\end{document}
